% SGA 3，第 XI 篇；根据英文重建源翻译。
% 范围：第 4 节全部，从标题、定理 4.1 至注 4.6；止于第 5 节之前。
% 权威来源：Polo--Gille Expo11.pdf，组件局部页 20--25；合卷页 742--747；印刷／原文页边页 157--164。
% Jacob Reinhold 的 CC BY 4.0 Markdown 谱系仅用于措辞比较；文字、公式、编号、注释及图均以 Polo--Gille PDF 为准。
% 精确续接／终点：第 5 节，Expo11 局部页 25／合卷页 747／原文页边页 165。

\SGARefTarget{sga3:xi:section:4}\section*{4. 平滑仿射群的乘法型\\子群概形}\addcontentsline{toc}{section}{4. 平滑仿射群的乘法型子群概形}

本篇的主要结果是 \SGARefLink{sga3:xi:theorem:4:1}{下述结果}。

\medskip\noindent\SGARefTarget{sga3:xi:theorem:4:1}\textbf{定理 4.1。} --- \textit{设 \(S\) 为预概形，\(G\) 为在 \(S\) 上平滑且仿射的 \(S\)-群预概形，设} \[
  F\colon(\mathbf{Sch})^\circ/S\longrightarrow(\mathbf{Ens})
\]\textit{为由下式定义的函子：} \SGARefTarget{sga3:xi:display:4:1:functor} \[
  F(T)=\{\text{\(G_T\) 的乘法型子群}\}.
\]\textit{则 \(F\) 可表，且在 \(S\) 上平滑并分离。}

立即记下 \SGARefLink{sga3:xi:corollary:4:2}{下述变体}。

\medskip\noindent\SGARefTarget{sga3:xi:corollary:4:2}\textbf{推论 4.2。} --- \textit{设 \(G\)、\(H\) 为两个 \(S\)-群预概形，其中 \(G\) 在 \(S\) 上平滑且仿射，\(H\) 为有限型乘法型群。则 \(\operatorname{Hom}_{S\text{-}\mathrm{gr}}(H,G)\) 可表，且在 \(S\) 上平滑并分离。}

事实上，当 \(H\) 在 \(S\) 上平滑时，\(H\times_S G\) 亦然，可以对后者应用 \SGARefLink{sga3:xi:theorem:4:1}{4.1}。通过考虑图子群，\(\operatorname{Hom}_{S\text{-}\mathrm{gr}}(H,G)\) 成为 \(H\times_S G\) 的乘法型子群函子 \(F\) 的子函子 \(F'\)：具体地，\(F'(T)\) 为 \[
  (H\times_S G)_T=H_T\times_TG_T
\] 中使第一投影所诱导的同态 \SGARefTarget{sga3:xi:display:4:2:first-projection} \[
  K\longrightarrow H_T,
\] 为同构的乘法型子群 \(K\) 的集合。由 \SGARefLink{sga3:ix:corollary:2:9}{第 IX 篇，2.9}，典范态射 \(F'\to F\) 为开浸入。由于 \(F\) 可表，\(F'\) 也可表，并由 \(F\) 的开子预概形表示；由于 \(F\) 在 \(S\) 上分离且平滑，\(F'\) 亦然。

当 \(H\) 在 \(S\) 上有限时，只需应用 \SGARefLink{sga3:xi:proposition:3:12:item:b}{3.12(b)} 即得 \(\operatorname{Hom}_{S\text{-}\mathrm{gr}}(H,G)\) 的可表性。当 \(H=H_1\times_SH_2\)，且 \(H_1\) 在 \(S\) 上平滑、\(H_2\) 在 \(S\) 上有限时，\(\operatorname{Hom}_{T\text{-}\mathrm{gr}}(H_T,G_T)\) 与 \SGARefTarget{sga3:xi:display:4:2:product-hom} \[
 \operatorname{Hom}_{T\text{-}\mathrm{gr}}\bigl((H_1)_T,G_T\bigr)
 \times
 \operatorname{Hom}_{T\text{-}\mathrm{gr}}\bigl((H_2)_T,G_T\bigr)
\] 中由满足 \(u_1\mathbin{\cdot}u_2=u_2\mathbin{\cdot}u_1\) 的对 \((u_1,u_2)\) 构成的子集等同。因此，\(\operatorname{Hom}_{S\text{-}\mathrm{gr}}(H,G)\) 由 \SGARefTarget{sga3:xi:display:4:2:representing-product} \[
 \operatorname{Hom}_{S\text{-}\mathrm{gr}}(H_1,G)
 \times_S
 \operatorname{Hom}_{S\text{-}\mathrm{gr}}(H_2,G)=X.
\] 的闭子预概形表示。事实上，应用 \SGARefLink{sga3:viii:examples:6:5:item:b:01}{第 VIII 篇，6.5(b)}，取 \(Y=H\)、\(Z=G\)，且 \(q_1\)、\(q_2\) 分别由 \[
  (u_1,u_2)\longmapsto u_1\mathbin{\cdot}u_2
  \qquad\text{且}\qquad
  (u_1,u_2)\longmapsto u_2\mathbin{\cdot}u_1.
\] 定义。

一般情形中，问题在 \(S\) 上对 Zariski 拓扑是局部的，因此可以假设 \(S\) 仿射，\(H\) 在 \(S\) 上具有常类型。由于 \(H\) 拟等平凡（\SGARefLink{sga3:x:corollary:4-5}{第 X 篇，4.5}），存在满射平展态射 \(S'\to S\)，其中 \(S'\) 仿射，使 \(H\) 分裂，即 \(H'=H_{S'}\) 可对角化。于是 \SGARefLink{sga3:xi:corollary:4:2}{上述结果} 适用，因为可对角化群是可对角化环面与有限可对角化群的乘积。还须说明 \(S'\)-预概形 \[
  \operatorname{Hom}_{S'\text{-}\mathrm{gr}}(H',G')=X'
\] 上所得的下降数据有效。这由 \SGARefLink{sga3:x:lemma:5-4}{第 X 篇，5.4} 的论证得到，只需注意那里假设 \(X'\to S'\) 分离且局部拟有限，仅用于保证 \(X'\) 的每个在 \(S'\) 上拟紧的开部分在 \(S'\) 上拟仿射。这里此性质仍成立，因为 \SGARefLink{sga3:xi:theorem:4:1}{4.1} 的函子 \(F\) 具有这一性质（将在 \SGARefLink{sga3:xi:theorem:4:1}{4.1} 的证明中看到），由此容易推出。此论证，或这种简短分解的任何类似变体，建立了 \(\operatorname{Hom}_{S\text{-}\mathrm{gr}}(H,G)\) 的可表性，同时表明它在 \(S\) 上分离。它在 \(S\) 上局部有限呈示，例如可由 \SGARefLink{sga3:xi:remark:3:2}{3.2} 中已指出的此函子“与环的滤过直极限交换”这一事实推出。最后，由 \SGARefLink{sga3:xi:proposition:2:1}{2.1}，此函子在 \(S\) 上形式平滑，因而在 \(S\) 上平滑。

\medskip\noindent\SGARefTarget{sga3:xi:remark:4:3}\textbf{注 4.3。} --- 我们从 \SGARefLink{sga3:xi:theorem:4:1}{4.1} 推出了 \SGARefLink{sga3:xi:corollary:4:2}{4.2}；只有当 \(H\) 也在 \(S\) 上平滑时，这才是直接的。要在一般情形中不作曲折论证而得出此推论，需要 \SGARefLink{sga3:xi:theorem:4:1}{4.1} 中的可表性断言不假设 \(G\) 在 \(S\) 上平滑，只假设它在 \(S\) 上仿射且有限呈示。当然，此时 \(F\) 一般不再在 \(S\) 上平滑。\SGARefLink{sga3:xi:theorem:4:1}{4.1} 在这些更一般的假设下仍成立，这一点几乎没有疑问，但证明看来更为细致，因为不能再使用 \SGARefLink{sga3:xi:corollary:3:8}{3.8}。\footnote{当 \(G\) 在 \(S\) 上平坦且拟仿射，并具有连通纤维时，只要限制于 \(G\) 的中心子环面，确已证明此事（\SGARefLink{sga3:xv:theorem:8-8}{第 XV 篇，8.8}）。}\footnote{编者注：当 \(H\) 在正规局部 Noether 基 \(S\) 上平滑（未必仿射）时，M.~Raynaud 证明，\(H\) 的子环面函子的最大可表开子函子，是若干在 \(S\) 上平滑且仿射的开集的不交和。这是\emph{群概形与齐性空间上的丰沛层〔Faisceaux amples sur les schémas en groupes et sur les espaces homogènes〕}，Lecture Notes in Mathematics 119 (1970) 中的定理 IX.9.26。}

不过应指出，当 \(G\) 为 \(S\) 上仿射平滑群 \(G_0\) 的闭子群时，\(G\) 的乘法型子群函子 \(F\) 可由表示 \(G_0\) 的类似函子 \(F_0\) 的预概形的闭子预概形表示，而后者由 \SGARefLink{sga3:xi:theorem:4:1}{4.1} 涵盖；应用 \SGARefLink{sga3:viii:theorem:6:4}{第 VIII 篇，6.4} 容易得到这一点。这也提出了如下问题：在 \(S\) 上仿射且有限呈示的仿射 \(S\)-群概形 \(G\)，是否同构于某个适当 \(n\) 对应的 \(\operatorname{GL}(n)_S\) 的子群概形。当 \(S\) 为域的谱时，这成立（参见第 \(\mathrm{VI}_{\mathrm B}\) 讲，11.11）；可惜一般并不成立，甚至对环面也不成立；见 \SGARefLink{sga3:xi:remark:4:6}{4.6}。最后，也可完全用证明 \SGARefLink{sga3:xi:theorem:4:1}{4.1} 的同一方法直接证明 \SGARefLink{sga3:xi:corollary:4:2}{4.2}。

现在证明 \SGARefLink{sga3:xi:theorem:4:1}{4.1}。由于函子 \(F\) 显然具有局部性质，可以假设 \(S\) 仿射，写成 \(S=\operatorname{Spec}(A)\)，其中 \(A\) 为环。将 \(A\) 视为其在 \(\mathbb Z\) 上有限型的子环的滤过直极限，并注意 \(G\) 下降自某个此类子环上的平滑仿射群（EGA IV，8），便归约到 \(S\) 为 Noether 概形的情形。

对每个整数 \(n>0\)，设 \(T_n\) 为与 \(F\) 同样定义、但仅考虑 \(G_T\) 中满足 \(n\cdot\operatorname{id}_H=0\)，即满足 \({}_nH=H\) 的乘法型子群 \(H\) 的函子。按整除关系对正整数集 \(I\) 排序。当 \(m\) 为 \(n\) 的倍数时，用 \[
  p_{n,m}(H)={}_nH=\operatorname{Ker}
  \bigl(n\cdot\operatorname{id}_H\bigr).
\] 定义 \SGARefTarget{sga3:xi:display:4:1:p-n-m} \[
  p_{n,m}\colon T_m\longrightarrow T_n
\]。这样，各 \(T_n\) 构成 \(S\) 上的函子逆系。由 \SGARefLink{sga3:xi:proposition:3:12:item:a}{3.12(a)}，各函子 \(T_n\) 可表，且在 \(S\) 上仿射并有限型。类似地，用 \(u_n(H)={}_nH\) 定义态射 \SGARefTarget{sga3:xi:display:4:1:u-n} \[
  u_n\colon F\longrightarrow T_n
\]。于是得到态射 \SGARefTarget{sga3:xi:display:4:1:u-inverse-limit} \[
  u\colon F\longrightarrow T=\varprojlim_n T_n,
\]，其中逆极限在 \(S\) 上的函子范畴中取。由于各 \(T_n\) 可表且在 \(S\) 上仿射，\(T\) 亦然：它是定义各 \(T_n\) 的 \(S\) 上拟凝聚代数的滤过直极限的谱。当然，\(T\) 一般不在 \(S\) 上有限型。

我们将应用 \SGARefLink{sga3:xi:proposition:3:7}{3.7}。只需验证其中的条件 \textit{\SGARefLink{sga3:xi:proposition:3:7:item:a}{(a)}}--\textit{\SGARefLink{sga3:xi:proposition:3:7:item:d}{(d)}}；这些条件将蕴含 \(F\) 可由 \(S\) 上局部有限型的预概形表示。进而由 \SGARefLink{sga3:xi:proposition:2:1:bis}{2.1 bis} 得知 \(F\) 甚至在 \(S\) 上平滑。此外，由于 \(F\) 为 \(T\) 的子函子，而后者在 \(S\) 上仿射，故 \(F\) 在 \(S\) 上分离：它在 \(T\) 上分离，而 \(T\) 在 \(S\) 上分离。还可以立即证明上面使用的附加断言：每个在 \(S\) 上拟紧的开集 \(U\subset F\) 在 \(S\) 上拟仿射。这由 \SGARefLink{sga3:xi:corollary:3:11}{3.11} 及各 \(T_n\) 在 \(S\) 上仿射得到。

现在验证 \SGARefLink{sga3:xi:proposition:3:7}{3.7} 的各个条件。

\smallskip\noindent\SGARefTarget{sga3:xi:theorem:4:1:condition:a}\textit{(a)} 由 SGA~1，第八讲的下降理论，\(F\) 为忠实平坦拟紧拓扑的层；这里此理论适用，因为 \(T\) 上的乘法型群在 \(T\) 上仿射。\footnote{印刷原文此处写“在 \(S\) 上仿射”；下降断言涉及 \(T\) 上的群，因此所需的基是 \(T\)。}由 EGA IV，8 的标准工具，它与环的滤过直极限交换。现说明它与局部 Artin 环的理想进逆极限交换。对 \SGARefLink{sga3:xi:corollary:4:2}{4.2} 中考虑的类似函子，这恰是 \SGARefLink{sga3:ix:theorem:7:1}{第 IX 篇，7.1} 在 \(A\) 为赋予通常拓扑的定义理想的完备 Noether 局部环时的特例。注意，正是在这里，本质地使用了 \(G\) 仿射的假设。在当前情形中，我们归约到 \SGARefLink{sga3:xi:lemma:4:4}{下述引理}。

\medskip\noindent\SGARefTarget{sga3:xi:lemma:4:4}\textbf{引理 4.4。} --- \textit{设 \(A\) 为完备 Noether 局部环，极大理想为 \(\mathfrak m\)，\(G\) 为 \(S=\operatorname{Spec}(A)\) 上的仿射群概形，对每个整数 \(n\geq0\)，置} \SGARefTarget{sga3:xi:display:4:4:adic-thickenings} \[
  S_n=\operatorname{Spec}\bigl(A/\mathfrak m^{\,n+1}\bigr),
  \qquad
  G_n=G\times_SS_n.
\]\textit{。对每个 \(n\)，设 \(H_n\) 为 \(G_n\) 的有限型乘法型子群，且假设当 \(m\geq n\) 时，\(H_n\) 由 \(H_m\) 经约化得到。则存在唯一的 \(G\) 的乘法型子群 \(H\)，其各个约化为各 \(H_n\)。}

由 \SGARefLink{sga3:x:theorem:3:2}{第 X 篇，3.2}，存在 \(S\) 上的乘法型群 \(H\)，它必然有限型且等平凡，在唯一同构意义下确定，并带有同构 \[
  H\times_SS_0\simeq H_0.
\]。这里使用了 \(H_0\) 为域上的有限型群、因而等平凡这一事实；见 \SGARefLink{sga3:x:proposition:1:4}{第 X 篇，1.4}。由 \SGARefLink{sga3:x:proposition:2:1}{第 X 篇，2.1}，对每个 \(n\)，同构 \(H\times_SS_0\simeq H_0\) 唯一地提升为同构 \(H\times_SS_n\simeq H_n\)。已引用的 \SGARefLink{sga3:ix:theorem:7:1}{第 IX 篇，7.1} 现表明，各同态 \(H_n\to G_n\) 来自唯一的 \(S\)-群同态 \SGARefTarget{sga3:xi:display:4:4:u} \[
  u\colon H\longrightarrow G.
\]。由 \SGARefLink{sga3:ix:corollary:6:6}{第 IX 篇，6.6}，\SGARefLink{sga3:xi:display:4:4:u}{此同态} 为单态射，因为 \(u_0\colon H_0\to G_0\) 是单态射。这就证明了 \SGARefLink{sga3:xi:lemma:4:4}{4.4}。

\smallskip\noindent\SGARefTarget{sga3:xi:theorem:4:1:condition:b}\textit{(b)} 态射 \(u\colon F\to T\) 为单态射。这由稠密性定理 \SGARefLink{sga3:ix:theorem:4:7}{第 IX 篇，4.7} 的形式 \SGARefLink{sga3:ix:corollary:4:8:item:b}{推论 4.8(b)} 推出。对于当前应用，关键在于这一结果在不必 Noether 的基上也成立。由于函子 \(T\) 不与环的滤过直极限交换，不能事先归约到 Noether 情形。

\smallskip\noindent\SGARefTarget{sga3:xi:theorem:4:1:condition:b-prime}\textit{(\(b'\))} 态射 \(u\colon F\to T\) 无穷小平展；换言之，有 \SGARefLink{sga3:xi:lemma:4:5}{下述引理}。

\medskip\noindent\SGARefTarget{sga3:xi:lemma:4:5}\textbf{引理 4.5。} --- \textit{设 \(A\) 为剩余域为 \(k\) 的 Artin 局部环，\(S=\operatorname{Spec}(A)\)，\(I\subset\operatorname{rad}(A)\) 为理想，置 \(S_0=\operatorname{Spec}(A/I)\)。设 \(G\) 为在 \(S\) 上平滑的 \(S\)-群预概形，置 \(G_0=G\times_SS_0\)。设给定 \(G_0\) 的乘法型子群 \(H_0\)，并对每个整数 \(n>0\) 给定 \(G\) 的乘法型子群 \(H(n)\)，使得}\begin{enumerate}\item[\SGARefTarget{sga3:xi:lemma:4:5:item:1}\(1^\circ\)]\textit{对 \(n\) 的每个倍数 \(m\)，有 \(H(n)={}_nH(m)\)；并且}\item[\SGARefTarget{sga3:xi:lemma:4:5:item:2}\(2^\circ\)]\textit{\(H(n)_0={}_nH_0\)。}\end{enumerate}\textit{则存在唯一的 \(G\) 的乘法型子群 \(H\)，使得对每个 \(n\)，有 \({}_nH=H(n)\)。}

唯一性已包含在 \textit{\SGARefLink{sga3:xi:theorem:4:1:condition:b}{(b)}} 中。对存在性，直接归纳可归约到 \(I\mathfrak m=0\) 的情形，其中 \(\mathfrak m\) 为 \(A\) 的极大理想。\footnote{印刷原文此处突然写作 \(J\mathfrak m=0\)，并在随后的论证中一直使用 \(J\)，尽管陈述中只引入了 \(I\)。这里始终使用同一理想 \(I\)。}置 \[
  k=A/\mathfrak m,\qquad
  S_0=\operatorname{Spec}(k),\qquad
  G_0=G\times_SS_0,
\]，设 \(\mathfrak g_0\)、\(\mathfrak h_0\) 分别为 \(G_0\)、\(H_0\) 的 Lie 代数。有典范群同构 \SGARefTarget{sga3:xi:display:4:5:infinitesimal-kernel} \[
  \mathfrak g_0\otimes_kI
  \simeq
  \operatorname{Ker}\bigl(G(S)\longrightarrow G(S_0)\bigr);
\]，见第三讲。由 \SGARefLink{sga3:ix:theorem:3:6:bis}{第 IX 篇，3.6 bis} 和 \SGARefPending{sga3:xi:proposition:3:7}{3.7}，存在 \(G\) 的乘法型子群 \(H\)，其约化为 \(H_0\)；这样的 \(H\) 在由 \[
  N=\operatorname{Ker}\bigl(G(S)\longrightarrow G(S_0)\bigr).
\] 的元素给出的内自同构意义下确定。因此 \(H_0\) 的提升 \(H\) 的集合 \(P\) 为 \(N/M\) 作用下的主齐性集，其中 \[
  M=\{g\in N\mid \operatorname{int}(g)\mathbin{\cdot}H=H\}.
\]。

容易看出（第三讲），此子群恰为 \(\mathfrak g_0\otimes_kI\) 中在 \(H_0\) 下不变的向量子空间，其中 \(H_0\) 通过 \(G_0\) 的伴随表示所诱导的表示作用于 \(\mathfrak g_0\otimes_kI\)。同样，\(H(n)_0={}_nH_0\) 的提升集合 \(P(n)\) 是 \(N/M(n)\) 作用下的主齐性集，其中 \(M(n)\) 为 \(\mathfrak g_0\otimes_kI=N\) 中在 \({}_nH_0\) 下不变的向量子空间。由稠密性定理 \SGARefLink{sga3:ix:theorem:4:7}{第 IX 篇，4.7} 容易看出，对充分大的 \(n\) 有 \(M=M(n)\)，其中“大”相对于正整数的整除偏序而言。因此，从 \(P\) 到 \(P(n)\) 的自然映射 \[
  H\longmapsto{}_nH
\] 对充分大的 \(n\) 为双射；它与 \(N\) 的作用相容，从而与作用群之间的同态 \[
  N/M\longrightarrow N/M(n)
\] 相容。\SGARefLink{sga3:xi:lemma:4:5}{4.5} 的结论立即随之成立。

为验证 \SGARefLink{sga3:xi:proposition:3:7}{3.7} 的条件 \textit{\SGARefLink{sga3:xi:proposition:3:7:item:c}{(c)}} 和 \textit{\SGARefLink{sga3:xi:proposition:3:7:item:d}{(d)}}，使用 \SGARefLink{sga3:xi:corollary:3:8}{3.8}，它使问题归约为验证下面的 \textit{(\(c'\))} 和 \textit{(\(d'\))}。

\smallskip\noindent\SGARefTarget{sga3:xi:theorem:4:1:condition:c-prime}\textit{(\(c'\))} 各 \(T_n\) 在 \(S\) 上平滑，且过渡态射 \(p_{n,m}\colon T_m\to T_n\) 平滑。

这恰是 \SGARefLink{sga3:xi:corollary:2:2:bis}{2.2 bis}。

\smallskip\noindent\SGARefTarget{sga3:xi:theorem:4:1:condition:d-prime}\textit{(\(d'\))} 用 \SGARefLink{sga3:xi:corollary:3:8}{3.8} 的记号，\(F\) 的取值于 \(S\) 上的域 \(k\) 的点 \(\xi\)，就是 \(G_0=G_k\) 的乘法型子群 \(H_0\)。重复上面的 \SGARefLink{sga3:xi:theorem:4:1:condition:b-prime}{\textit{(\(b'\))}} 中的论证可知，\SGARefLink{sga3:xi:corollary:3:8}{3.8} 所考虑的整数 \(d(\xi)\) 为 \SGARefTarget{sga3:xi:display:4:1:d-xi-quotient} \[
  d(\xi)=\dim_k\bigl(\mathfrak g_0/\mathfrak g_0^{H_0}\bigr),
\]，其中 \(\mathfrak g_0\) 为 \(G_0\) 的 Lie 代数，\(\mathfrak g_0^{H_0}\) 为在 \(H_0\) 下不变的向量子空间。这是有限整数，故 \SGARefLink{sga3:xi:corollary:3:8}{3.8} 的条件 \SGARefLink{sga3:xi:corollary:3:8:item:d-prime:1}{\(\textit{(d')}1^\circ\textit{)}\)} 成立。

用 \SGARefLink{sga3:xi:corollary:3:8:item:d-prime:2}{\(\textit{(d')}2^\circ\textit{)}\)} 的记号，给定态射 \(v\colon
X\to F\)，等于给定 \(G_X\) 的乘法型子群 \(H\)。对 \(x\in X\)，整数 \(d(\xi_x)\)\footnote{印刷原文此处仅给出 \(\bigl(\mathfrak g\otimes\kappa(x)\bigr)^{H_x}\)，后面又写 \(\operatorname{rank}_x\mathfrak g^H\)，均遗漏了商。这一读法同时与 \SGARefLink{sga3:xi:display:4:1:d-xi-quotient}{紧前的公式} \(d(\xi)=\dim(\mathfrak g_0/\mathfrak g_0^{H_0})\) 及提升挠子 \(N/M\) 矛盾。两处均恢复了商。}即 \SGARefTarget{sga3:xi:display:4:1:d-xi-x-invariants} \[
  \bigl(\mathfrak g\otimes\kappa(x)\bigr)
  \big/
  \bigl(\mathfrak g\otimes\kappa(x)\bigr)^{H_x}.
\] 的维数，其中 \(\mathfrak g\) 为 \(G_X\) 的 Lie 代数层。由于 \(G_X\) 在 \(X\) 上平滑，它在 \(X\) 上局部自由且有限型，而 \(\mathfrak g\otimes\kappa(x)\) 为 \(G_X\) 在 \(x\) 处的纤维 \(G_x\) 的 Lie 代数；这里 \(H_x\) 为 \(H\) 在 \(x\) 处的纤维。而 \(\mathfrak g\) 为 \(X\) 上的局部自由模\footnote{印刷原文此处写“在 \(S\) 上”，但 \(\mathfrak g\) 为 \(G_X\) 的 Lie 代数层；因此基应为 \(X\)。}，乘法型群 \(H\) 作用于其上，所以不变子函子 \(\mathfrak g^H\) 由 \(\mathfrak g\) 中局部为直和项、因而局部自由的子层给出。通过下降，可归约到 \(H\) 可对角化的情形，然后应用 \SGARefLink{sga3:i:proposition:4-7-3}{第 I 篇，4.7.3}，只需注意不变子层对应于零次分量。因此 \[
  d(\xi_x)=\operatorname{rank}_x
  \bigl(\mathfrak g/\mathfrak g^H\bigr),
\]，所以这是 \(x\) 的局部常值函数。这就证明了条件 \SGARefLink{sga3:xi:theorem:4:1:condition:d-prime}{\textit{(\(d'\))}}。

这样就验证了 \SGARefLink{sga3:xi:proposition:3:7}{3.7} 的各条件，完成了 \SGARefLink{sga3:xi:theorem:4:1}{4.1} 的证明。

\medskip\noindent\SGARefTarget{sga3:xi:remark:4:6}\textbf{注 4.6。} --- 当 \(G=\operatorname{GL}(n)_S\) 时，利用 \SGARefLink{sga3:i:proposition:4-7-3}{第 I 篇，4.7.3} 可以给出 \SGARefLink{sga3:xi:theorem:4:1}{4.1} 的简单得多且更明确的直接证明。该证明还表明，此时模概形为仿射 \(S\)-概形的不交和。按 \SGARefLink{sga3:xi:remark:4:3}{4.3} 所指方式进行，只要 \(G\) 为形如 \(\operatorname{GL}(n)_S\) 的群的闭子群，就得到相同的结果。

但不能据此认为，表示 \SGARefLink{sga3:xi:theorem:4:1}{4.1} 或 \SGARefLink{sga3:xi:corollary:4:2}{4.2} 中函子的预概形总是仿射 \(S\)-概形的不交和。例如，设 \(H\) 为局部 Noether 预概形 \(S\) 上的有限型乘法型群。由 \SGARefLink{sga3:x:proposition:5-11}{第 X 篇，5.11}，\(H\) 等平凡当且仅当 \[
  \operatorname{Hom}_{S\text{-}\mathrm{gr}}(H,\mathbb G_m)
\] 为 \(S\) 上仿射预概形的不交和。但如 \SGARefLink{sga3:x:paragraph:1:6}{第 X 篇，1.6} 所指出，可能存在相对维数为 \(2\) 的非等平凡环面 \(H\)。对此类环面，\(\operatorname{Hom}_{S\text{-}\mathrm{gr}}(H,\mathbb G_m)\) 不是在 \(S\) 上仿射的 \(S\)-预概形的不交和，类似地，“对偶”的扭常群 \[
  \operatorname{Hom}_{S\text{-}\mathrm{gr}}(\mathbb G_m,H)
\] 也不是这样的和。事实上，若两个有限呈示的交换扭常群 \(R,R'\) 互为对偶，\[
  R'=\operatorname{Hom}_{S\text{-}\mathrm{gr}}(R,\mathbb Z_S),
\]，则容易看出其中一个等平凡当且仅当另一个等平凡。

最后一点也表明，这样的 \(H\) 不同构于形如 \(\operatorname{GL}(n)_S\) 的群的子群。更准确地说，可以证明：连通局部 Noether 基 \(S\) 上的有限型乘法型群，同构于形如 \(\operatorname{Aut}_{\mathrm{Mod}}(E)\) 的群的子群（其中 \(E\) 为 \(S\) 上有限型的局部自由模），当且仅当它等平凡。最后，在前两个例子中取 \(G=H\times\mathbb G_m\)，便得到一个例子，使表示 \SGARefLink{sga3:xi:theorem:4:1}{4.1} 中函子 \(F\) 的预概形不是仿射 \(S\)-预概形的不交和；若愿意，可取 \(G\) 为相对维数 \(3\) 的环面。


